% Fuente: LAFROX-Subdocumento8.md
% Texto y datos conservados; formato adaptado a lafrox.cls.

\section*{Introducción a los Riesgos}
\addcontentsline{toc}{section}{Introducción a los Riesgos}

La solución de LafroX integra recepción y trazabilidad, inventario, frío, preventa, reparto y cobranza en E1, y canal moderno, portales y costo de servir en E2. Sus riesgos dependen de una operación que despacha 96 camiones entre 05:30 y 07:00, mantiene los CD 24 horas sin WAN y el terreno 14 horas sin señal, conserva ERP como único emisor tributario y despliega sin detener rutas (Distribuidora Puelche S.A., 2026c, Cap. 10). Este plan conecta las decisiones comerciales del SD2, alcance del SD3, contratos y dimensionamiento del SD4, gobierno del SD6 y programación del SD7 (LafroX, 2026).

El cuerpo establece método y decisiones; los Anexos 8.A–8.F contienen registro ampliado, FMEA, escenarios, reservas y condiciones de evidencia; el T-16 conserva el formato obligatorio separado (Distribuidora Puelche S.A., 2026d, sección 11). Las condiciones que deben cerrarse antes de cada hito, incluida la verificación del RPO y del riesgo residual que el SD4 declara en la sección 4.3.2.4, se registran en el Anexo 8.E con responsable y plazo.

\section{Plan de riesgos}

Esta sección define cómo se gestionan los riesgos durante los 56 meses: el ciclo y sus instancias, los roles y su autoridad, las escalas con que se analizan y las reglas de registro y cierre.

\subsection{Enfoque y ciclo}

El registro vivo sigue identificación, análisis, respuesta y seguimiento conforme a la norma ISO 31000 (International Organization for Standardization [ISO], 2018), que exige el RT-19.04 de las Bases Técnicas Transversales (Distribuidora Puelche S.A., 2026b). Cada entrada distingue causa, evento incierto y consecuencia, sin convertir una inconsistencia observada en riesgo futuro. El Anexo 8.E separa problemas y dependencias actuales.

JP mantiene el registro único. Los líderes identifican riesgos al revisar interfaces, pruebas, capacidad y despliegues; verifican disparadores semanalmente y ante cada cambio. El Comité de Proyecto quincenal revisa todos los riesgos abiertos, responsables, plazos y evidencia. Los comités Ejecutivo y de Arquitectura mensuales deciden escalamiento y cambios de sus ámbitos, y desde el mes 13 el Comité de Operación mensual revisa los riesgos de servicio: niveles de atención, incidentes, capacidad y continuidad (SD6, sección 6.1.6). Una amenaza a continuidad, seguridad o hito se escala inmediatamente. Marcha blanca exige seguimiento diario; operación mantiene vigilancia diaria de incidentes y colas y revisión en comités. En cada comité se compara la reserva restante con el riesgo remanente (análisis de reserva), y al cierre de cada etapa, en el H5 y el H10, una auditoría de riesgos comprueba que el proceso funciona (PMI, 2017, p. 456). El Comité Ejecutivo recibe cada mes un extracto con los diez riesgos de mayor exposición; el registro completo queda en el repositorio del proyecto.

\subsection{Roles y autoridad}

La Tabla 8.1 asigna a cada líder del proyecto el ámbito de riesgos que vigila en el registro.


\renewcommand{\thetable}{8.1}
\begin{tablalafrox}{Ámbito de riesgos por rol. Fuente: elaboración propia a partir del Capítulo 1, Tabla 1.3, y del Formulario T-16}{tab:sd8-cuerpo-1}{>{\hsize=0.75287\hsize\linewidth=\hsize\RaggedRight\arraybackslash}X >{\hsize=1.24713\hsize\linewidth=\hsize\RaggedRight\arraybackslash}X}{\cab{Rol} & \cab{Ámbito de riesgos que vigila}}
JP — Alex Aravena & Registro, fechas, contraparte, recursos y escalamiento \\
ARQ — Bastián Trejo & Interfaces, portabilidad, continuidad y tratamiento del riesgo residual de RPO \\
SEG — Álvaro Catalán & Seguridad, acceso, datos y respuesta a incidentes \\
DAT — Leandro Chamorro & Migración, conciliación, trazabilidad y modelos \\
DES — Tomás Pérez & Módulos, integración, idempotencia y equipos separados E1/E2 \\
CAL — Maximiliano Miño & Pruebas, criterios y evidencia de cierre \\
SRE — Guillermo Castillo & Capacidad, infraestructura, enlaces, continuidad y atención \\
IMP — Patricio Henríquez & Olas, usuarios, formación, acuerdos y adopción \\
\end{tablalafrox}

Dirigen equipos; no ejecutan solos todas las HH. La Contraparte Técnica del CLIENTE acepta entregables mediante actas; Operaciones autoriza cortes y continuidad de su ámbito; Calidad del CLIENTE valida decisiones sanitarias. Su disponibilidad se acuerda en la agenda de decisiones del mes 1 y se sigue con R8-12. Los cambios físicos siguen el gobierno de arquitectura, sin modificar silenciosamente SD4/T-11.

\subsection{Escalas previas al análisis}

P/I/D se asignan como juicios ordinales iniciales sustentados en exposición y controles descritos, no como frecuencias medidas; para el análisis cuantitativo, P e I se calibran con los tramos de la Tabla 8.3. La Tabla 8.2 define los cinco niveles de cada escala. Horizonte: hasta entregar el control de cada ficha y, para riesgos recurrentes, los 56 meses del contrato. Cambiar el horizonte exige reevaluación.


\renewcommand{\thetable}{8.2}
\begin{tablalafrox}{Escalas ordinales de probabilidad, impacto y detección. Fuente: elaboración propia a partir de ISO (2018) e IEC (2018)}{tab:sd8-cuerpo-2}{>{\hsize=0.23658\hsize\linewidth=\hsize\RaggedRight\arraybackslash}X >{\hsize=1.21097\hsize\linewidth=\hsize\RaggedRight\arraybackslash}X >{\hsize=1.27842\hsize\linewidth=\hsize\RaggedRight\arraybackslash}X >{\hsize=1.27403\hsize\linewidth=\hsize\RaggedRight\arraybackslash}X}{\cab{Valor} & \cab{Probabilidad ordinal P} & \cab{Impacto I: mayor efecto aplicable} & \cab{Detección D: mayor valor, peor detección}}
1 & Excepcional, sin dependencia expuesta identificada & Corrección local sin afectar aceptación ni servicio & Control automático probado detecta antes de confirmar \\
2 & Posible con exposición limitada & Retrabajo absorbible en capacidad asignada & Ensayo sistemático detecta antes de habilitar \\
3 & Plausible por una dependencia expuesta & Consume contingencia o afecta servicio auxiliar & Requiere conciliación o ensayo específico \\
4 & Exposición recurrente o varias dependencias sin control demostrado & Amenaza hito, SLA o alcance obligatorio & Se aprecia al afectar operación o en revisión tardía \\
5 & Condiciones muy favorables al evento sin barrera eficaz & Detiene despacho crítico, compromete seguridad/datos/sanidad o impide aceptación & Sin evidencia de detección o reconocible después de pérdida \\
\end{tablalafrox}

Exposición E = P \ensuremath{\times} I (1–25): baja 1–3, moderada 4–7, alta 8–14 y crítica 15–25. FMEA agrega NPR = P \ensuremath{\times} I \ensuremath{\times} D (1–125), para ordenar dentro del nivel; desempates por impacto y proximidad del plazo. I = 5 exige escalamiento aunque E no alcance 15. No se asigna D = 1 a un control por describirlo.

El apetito de riesgo del proyecto se fija por nivel de exposición, con una regla de acción por zona. Un riesgo crítico (E de 15 a 25) o con I = 5 se evita, se mitiga o se escala antes de su hito, con su control como paquete de la EDT, aunque su retorno en HH sea menor que 1; si el control falla, se cambia el diseño. Si su causa no admite un control previo, como la demanda real de la mesa (R8-22), se acepta activamente con la contingencia ya dimensionada y su disparador. Uno alto (8 a 14) se mitiga con un control de la EDT y revisión quincenal en el Comité de Proyecto. Uno moderado (4 a 7) se acepta activamente, con reserva y disparador. Uno bajo (1 a 3) se acepta pasivamente, con revisión mensual. La tolerancia es cero días de atraso en los hitos del Art. 17° y en el despacho de 05:30 a 07:00, y ningún riesgo que afecte un requisito obligatorio se acepta sin tratamiento. El umbral de escalamiento es cualquier riesgo crítico o con I = 5, o un consumo de la reserva de contingencia mayor que el previsto para el período en la Tabla 8.6.

Para el análisis cuantitativo, cada valor de P se calibra con un tramo de probabilidad y cada valor de I con la fracción del esfuerzo de los paquetes afectados que el evento agregaría como retrabajo o trabajo adicional. Los paquetes afectados son los que cada ficha del Anexo 8.A nombra en «Fuente y EDT», con sus HH del Formulario T-15. La calibración es un juicio del equipo, no una frecuencia medida, y se revisa con los datos del proyecto. La Tabla 8.3 fija los valores usados.


\renewcommand{\thetable}{8.3}
\begin{tablalafrox}{Calibración cuantitativa de las escalas. Fuente: elaboración propia a partir del Formulario T-15}{tab:sd8-cuerpo-3}{>{\hsize=0.41951\hsize\linewidth=\hsize\RaggedRight\arraybackslash}X >{\hsize=1.41641\hsize\linewidth=\hsize\RaggedRight\arraybackslash}X >{\hsize=1.16408\hsize\linewidth=\hsize\RaggedRight\arraybackslash}X}{\cab{Valor} & \cab{Probabilidad (tramo y valor usado)} & \cab{Impacto (fracción del esfuerzo afectado)}}
1 & Menos de 10 \%; 5 \% & Sin esfuerzo adicional \\
2 & 10 a 30 \%; 20 \% & 5 \% \\
3 & 30 a 50 \%; 40 \% & 10 \% \\
4 & 50 a 70 \%; 60 \% & 20 \% \\
5 & Más de 70 \%; 80 \% & 30 \% \\
\end{tablalafrox}

El valor esperado de un riesgo es su probabilidad por su impacto en HH. Las horas sirven de unidad porque la oferta técnica no contiene montos (Distribuidora Puelche S.A., 2026a, Art. 50.2); la Oferta Económica valoriza esas horas.

\subsection{Registro y cierre}

El responsable actualiza fuente, estado, disparador, tratamiento, consumo de HH y evidencia. Estados: identificado, en tratamiento, control verificado o materializado. Materializar el evento abre un problema y conserva su ID. CAL verifica cierre técnico; JP registra la decisión; la aceptación contractual corresponde al CLIENTE. La puntuación residual se estima después de comprobar controles, no por asignar una mitigación. Ninguna aceptación de riesgo exime requisitos obligatorios.

\section{Identificación y Análisis de Riesgos}

Esta sección identifica los riesgos con una RBS y los analiza en dos niveles: uno cualitativo con FMEA y otro cuantitativo con el valor esperado de cada riesgo en horas hombre y una simulación de Monte Carlo sobre el cronograma por actividad del SD7.

\subsection{RBS y análisis separados}

La estructura de desglose de riesgos (RBS) tiene raíz «Riesgos de la propuesta LafroX» y tres ramas, una por cada análisis que exige el Formulario T-22: riesgo de la solución, del desarrollo y de la implantación. Dentro de cada rama, los riesgos se agrupan en las categorías técnica, organizacional, de proyecto, de seguridad y de operación que les corresponden. Los ID son únicos aunque una causa afecte varias ramas. La Figura 8.1 presenta la estructura con los 32 riesgos del Anexo 8.A.

\parte{figuras/rbs}


La figura muestra que la solución concentra 13 de los 32 riesgos y 11 de los 22 de nivel crítico de la Tabla B.1, porque reúne la continuidad del despacho, la reserva de stock y la evidencia de frío. Los dos riesgos de seguridad (R8-07 y R8-24) pertenecen a la rama de la solución, porque sus fichas los originan en los portales, móviles, integraciones y trazas que la arquitectura expone. El desarrollo se concentra en la categoría de proyecto (dotación, capacidad protegida de la Etapa 1, certificación de cadenas y revisión del CLIENTE), y la implantación, en la organizacional (rotación, transportistas, sindicato y adopción de innovaciones).

Solución: integridad de stock/custodia, ERP/DTE, autonomía, capacidad, RPO, seguridad, obsolescencia y proveedores. Desarrollo: interfaces sin documentación, recursos, contrapartes, conocimiento de rutas, migración, certificación externa, certificación paralela a la revisión del CLIENTE y refuerzo subcontratado de calidad. Implantación: congelamientos, cuatro semanas completas, suministros, adopción, relevos y atención. Anexo 8.A desarrolla 32 amenazas; 8.F cubre las cinco innovaciones y una oportunidad de diagnóstico con INN-02. Las tres ramas permiten preparar los análisis separados exigidos por T-22.

Los siete riesgos que el Caso 02 espera encontrar en este plan (Distribuidora Puelche S.A., 2026c, Cap. 19) tienen ficha propia: la pérdida del conocimiento de rutas (R8-13), la objeción sindical y el rechazo de los conductores de terceros (R8-21), la rotación en bodega (R8-20) y la ausencia de documentación de interfaces (R8-10). La dependencia de un solo enlace en Concepción la resuelve el SD4 con tres caminos independientes, y R8-03 la verifica con la prueba de 24 horas sin WAN en cada centro de distribución. El peak de septiembre coincide con las semanas de cierre de la marcha blanca de la Etapa 2, y R8-18 lo trata con la habilitación completa en agosto y la prueba de carga a 1,5 veces el peak antes del H11.

\subsection{Resultados cualitativos y FMEA}

El análisis de modos de falla y sus efectos (FMEA) sigue la IEC 60812 (International Electrotechnical Commission [IEC], 2018): a la probabilidad y el impacto agrega la detección, y su producto da el número de prioridad del riesgo (NPR). Anexo 8.B presenta P/I/D, exposición y NPR. La prioridad inicial se concentra en continuidad, RPO, ERP/DTE, seguridad, capacidad y aceptación. T-16 compara exposición; NPR no reemplaza su columna Expos.

La Tabla 8.4 extrae del Anexo 8.B, Tabla B.1, los cinco riesgos con mayor número de prioridad (NPR); el registro completo de los 32 está en ese anexo.


\renewcommand{\thetable}{8.4}
\begin{tablalafrox}{Riesgos con mayor NPR. Fuente: Anexo 8.B, Tabla B.1}{tab:sd8-cuerpo-4}{>{\hsize=2.74513\hsize\linewidth=\hsize\RaggedRight\arraybackslash}X >{\hsize=0.44728\hsize\linewidth=\hsize\RaggedRight\arraybackslash}X >{\hsize=0.44728\hsize\linewidth=\hsize\RaggedRight\arraybackslash}X >{\hsize=0.44728\hsize\linewidth=\hsize\RaggedRight\arraybackslash}X >{\hsize=0.91302\hsize\linewidth=\hsize\RaggedRight\arraybackslash}X}{\cab{Riesgo} & \cab{P} & \cab{I} & \cab{D} & \cab{NPR = P \ensuremath{\times} I \ensuremath{\times} D}}
R8-01 ERP indisponible o guía invalidada & 4 & 5 & 4 & 80 \\
R8-07 Ataque o exposición de datos críticos & 4 & 5 & 4 & 80 \\
R8-27 INN-03 estima vida remanente insegura & 4 & 5 & 4 & 80 \\
R8-05 Pérdida del sitio supera RPO & 3 & 5 & 5 & 75 \\
R8-02 Doble reserva o custodia en la coordinación de reserva & 4 & 5 & 3 & 60 \\
\end{tablalafrox}

Los cinco pertenecen a la rama de la solución y todos tienen impacto 5. Lo que los separa es la detección: R8-05 tiene la probabilidad más baja, pero su D = 5 (sólo se aprecia al perder el sitio) lo deja en el cuarto lugar, por sobre riesgos más probables. Por eso sus respuestas privilegian controles que se prueban antes del corte (ensayo de caída del ERP, prueba de concurrencia y ejercicio de recuperación) en vez de controles que sólo reaccionan cuando el evento ocurre.

No se suman puntuaciones como probabilidad del proyecto. Las relaciones importan: ausencia de contraparte puede atrasar interfaces/cadenas y eliminar semanas de evidencia; migración deficiente puede invalidar trazabilidad y aceptación. Los mismos efectos o consumos no se contabilizan dos veces.

\subsection{Análisis cuantitativo}

El análisis cuantitativo tiene dos partes (PMI, 2017, pp. 433–434). La primera calcula el valor esperado de cada riesgo con la calibración de la sección 8.1.3. El registro suma 15.076 HH de valor esperado, cerca de 8 \% de las 190.366 HH base del T-15. Cinco riesgos concentran 67,7 \% del total: productividad o dotación inferior al modelo (R8-11), mesa que no alcanza los niveles de atención (R8-22), marcha blanca que no cumple las seis condiciones (R8-18), uso de la capacidad protegida de la Etapa 1 por la Etapa 2 (R8-14) y doble reserva de stock (R8-02). El Anexo 8.B, Tabla B.2, presenta el cálculo de cada riesgo y su porcentaje acumulado.

La segunda parte simula 5.000 veces la red de los 163 paquetes con entregable del Formulario T-15, programados por actividad. En cada iteración, la duración de cada paquete varía con la distribución PERT de su tríada, cada riesgo ocurre con su probabilidad y, si ocurre, alarga sus paquetes afectados en la fracción de su impacto. La Tabla 8.5 informa, para cada hito, la fecha límite de entrega, las fechas que se alcanzan en la mitad (P50) y en el 80 \% (P80) de las iteraciones y la probabilidad de entregar a tiempo. Los hitos H6, H7, H11 y H12 no se simulan: son el inicio de cada marcha blanca y cada paso a producción, con mes fijo del Art. 17°, y su riesgo se trata con R8-18 y la reserva de contingencia (Anexo 8.C, sección C.3).


\renewcommand{\thetable}{8.5}
\begin{tablalafrox}{Resultado de la simulación de Monte Carlo por hito. Fuente: elaboración propia a partir del Formulario T-15 y del Anexo 8.C}{tab:sd8-cuerpo-5}{>{\hsize=0.55217\hsize\linewidth=\hsize\RaggedRight\arraybackslash}X >{\hsize=1.22287\hsize\linewidth=\hsize\RaggedRight\arraybackslash}X >{\hsize=1.09774\hsize\linewidth=\hsize\RaggedRight\arraybackslash}X >{\hsize=1.09774\hsize\linewidth=\hsize\RaggedRight\arraybackslash}X >{\hsize=1.02949\hsize\linewidth=\hsize\RaggedRight\arraybackslash}X}{\cab{Hito} & \cab{Fecha límite de entrega} & \cab{P50} & \cab{P80} & \cab{P(entrega a tiempo)}}
H1 & 17-03-2027 & 09-03-2027 & 10-03-2027 & > 99,9 \% \\
H2 & 17-05-2027 & 06-04-2027 & 09-04-2027 & > 99,9 \% \\
H3 & 16-07-2027 & 30-06-2027 & 09-07-2027 & 98,6 \% \\
H4 & 16-11-2027 & 28-10-2027 & 03-11-2027 & > 99,9 \% \\
H5 & 17-01-2028 & 27-12-2027 & 05-01-2028 & 97,7 \% \\
H8 & 17-03-2028 & 06-03-2028 & 09-03-2028 & 99,3 \% \\
H9 & 16-06-2028 & 06-06-2028 & 14-06-2028 & 89,9 \% \\
H10 & 17-07-2028 & 30-06-2028 & 07-07-2028 & 97,7 \% \\
\end{tablalafrox}

La fecha P80 de cada hito queda antes de su fecha límite: el plan se compromete al percentil 80 y se ejecuta sobre las fechas programadas, que están cerca del P50. Para lograrlo, el SD7 dimensionó las reservas de cada hito con esta simulación, adelantó la sala técnica, los ambientes y la integración de la Etapa 2, y refuerza la calidad con evaluadores subcontratados durante las certificaciones. El H9 es el hito más expuesto (89,9 \%), y su probabilidad depende sobre todo de R8-14, R8-11 y R8-15 (Anexo 8.C, Tabla C.4).

T-15 programa 202.774 HH, con soporte puente de 9.336 HH (parte de las 16.664 HH del período de los meses 16 a 20 del SD7, Tabla 7.4) y reserva E1 de 3.072 HH ya incluidos; la mesa y el SOC siguen las posiciones del SD4 y el calendario real. El máximo mensual es 69 personas equivalentes en el mes 15, y la construcción de la Etapa 1 ocupa a 48 personas de desarrollo a la vez, las 48 que declara la división de desarrollo del SD1. SEG e IMP se completan con contratación o subcontratación antes de aprobar la línea base, con asignación nominal por subventana (T-15 §5.7; R8-11 y E8-01). Los niveles de atención de la mesa se miden por contacto desde la marcha blanca de la Etapa 1, y su dotación se recalibra antes del mes 21 (R8-22).

\section{Plan de Acción a Riesgos}

Esta sección fija las respuestas a los riesgos analizados, las reservas que las financian en horas y las condiciones que deben cumplirse antes de declarar factible la propuesta.

\subsection{Respuestas y costo-beneficio técnico}

Cada ficha 8.A fija responsable, estrategia, disparador, plazo, mitigación, contingencia, efecto esperado, riesgo residual, riesgo secundario y evidencia de cierre. La estrategia sigue las cinco del PMBOK para amenazas (PMI, 2017, pp. 442–443) y se elige según quién controla la causa. Se mitigan 28 riesgos, porque LafroX controla su causa. Se escalan dos, R8-12 y R8-17, porque dependen de decisiones del CLIENTE. Se evita R8-14 separando los equipos de las dos etapas. Se acepta activamente R8-22, con reserva y disparador, porque no tiene un control previo rentable. Ninguno se transfiere: contratar a un tercero no traslada la obligación final de LafroX con el CLIENTE. La atención auxiliar en contingencia no acredita capacidades obligatorias desactivadas.

Cada respuesta deja un riesgo residual. Con la probabilidad un nivel más baja tras verificar cada control, el valor esperado del registro baja de 15.076 a 10.803 HH. Las respuestas también crean riesgos secundarios, que el registro trata como riesgos propios: certificar en paralelo a la revisión del CLIENTE crea R8-31, reforzar la calidad con evaluadores subcontratados crea R8-32, capturar evidencia de terreno para INN-02 crea R8-24, y separar los equipos de las dos etapas presiona la dotación de R8-11.

El costo-beneficio sigue la regla del PMBOK: una respuesta se justifica si reduce el valor esperado más de lo que cuesta (PMI, 2017, pp. 442–443). El Anexo 8.C, Tabla C.5, calcula para los 22 riesgos críticos el ahorro esperado, que es la diferencia entre el valor esperado inicial y el residual, y lo divide por las HH del paquete de control del T-15. El retorno supera 1 en tres controles: la nivelación de recursos frente a la productividad (R8-11, 17,0) y frente al uso de la capacidad de la Etapa 1 (R8-14, 4,2), y el plan de olas con la certificación de usuarios frente a la marcha blanca (R8-18, 2,1). En los demás, el ahorro medido sólo en HH de retrabajo es menor que el costo del control, porque ese impacto no incluye la detención del despacho, la sanción sanitaria ni el atraso de un hito, que son las consecuencias que hacen crítico al riesgo. Para ellos rige la regla del nivel crítico de la sección 8.1.3, y la mayoría de sus controles son pruebas o actas que las Bases exigen. El cumplimiento obligatorio prevalece sobre el retorno. Valorización, tarifas y reservas monetarias corresponden exclusivamente a la Oferta Económica (BA Art. 50.2).

\subsection{Reservas y cronograma}

La reserva de contingencia cubre los riesgos identificados, forma parte de la línea base y se dimensiona con la suma de sus valores esperados: 15.076 HH, el 7,9 \% de las 190.366 HH base del T-15 (PMI, 2017, p. 202; p. 443). Se dimensiona con el valor esperado inicial y no con el residual, porque ningún control está verificado al ofertar; cada control verificado libera la diferencia, hasta 4.273 HH si se verifican todos (Anexo 8.D). De ellas, 1.853 HH corresponden a R8-02, R8-04 y R8-14 (461 + 384 + 1.008 HH, Tabla B.2), los riesgos de corrección de la Etapa 1 que la capacidad protegida de 3.072 HH de los meses 13 a 20, ya incluida en el T-15, puede absorber según la regla del Anexo 8.D. Esa capacidad no puede usarse antes del mes 13 ni para la Etapa 2, de modo que no cubre a los demás riesgos, y la contingencia adicional es de 15.076 \ensuremath{-} 1.853 = 13.223 HH. La Tabla 8.6 reparte la contingencia por período según los meses de los paquetes afectados por cada riesgo; ese reparto es el reflejo de la reserva en el flujo de caja, y su valorización está en la Oferta Económica (Art. 50.2).


\renewcommand{\thetable}{8.6}
\begin{tablalafrox}{Reparto de la reserva de contingencia por período. Fuente: elaboración propia a partir del Anexo 8.B, Tabla B.2}{tab:sd8-cuerpo-6}{>{\hsize=1.09483\hsize\linewidth=\hsize\RaggedRight\arraybackslash}X >{\hsize=0.90517\hsize\linewidth=\hsize\RaggedRight\arraybackslash}X}{\cab{Período} & \cab{Contingencia (HH)}}
Meses 1–6 & 1.802 \\
Meses 7–12 & 4.090 \\
Meses 13–15 & 2.359 \\
Meses 16–21 & 3.287 \\
Meses 22–33 & 3.486 \\
Meses 34–56 & 51 \\
\textbf{Total} & \textbf{15.076} \\
\end{tablalafrox}

Los valores de cada período se redondean a la hora; el total se calcula sin redondear. El 88 \% de la reserva (13.222 de 15.076 HH) se concentra entre los meses 7 y 33, donde coinciden la construcción, las dos marchas blancas y el primer año de operación; después del mes 33 queda sólo el remanente de los riesgos de operación (51 HH).

La reserva de cronograma son las reservas de cada hito del T-15, Tabla 5.2, dimensionadas para que la fecha P80 quede antes de la fecha límite (sección 8.2.3). La reserva de gestión cubre riesgos no identificados: no forma parte de la línea base, la autoriza el Comité Ejecutivo, usarla exige actualizar la línea base y su monto se define en la Oferta Económica (BA Art. 50.2). JP solicita el uso de cualquier reserva con causa, perfiles, ventana e impacto; ninguna reserva se presta entre etapas.

Cada uso registra un cargo único por evento/mes/perfil y remanente. Riesgos correlacionados comparten consumo real; E1 no presta su reserva a E2. Los recursos adicionales requieren actualizar T-15 y calendario, sin ampliar automáticamente los 56 meses. Los meses 21 y 22 separan 2.774,54 HH de cierre y estabilización de implementación de la operación.

\subsection{Factibilidad y aceptación}

El Anexo 8.E registra las condiciones de evidencia: productividad/dotación, secuencias diarias y plazos de revisión/subsanación del Art.18.3, fecha contractual, continuidad de 96 despachos, RPO remoto, atención y evidencia láctea. El dimensionamiento del Capítulo 4 (Anexo 4-I y sección 4-W.5) ya incluye la coordinación de reserva de INT-03/04, con cuatro mensajes por línea de pedido; la proporción de líneas repartidas entre lotes o sitios y el drenaje de la cola se verifican en la prueba de concurrencia previa al H4. El RPO \ensuremath{\leq} 15 min se cumple con la replicación por los tres caminos independientes de Talca (SD4, sección 4.2.5.1). La falla simultánea de los tres caminos seguida de la destrucción del sitio queda como riesgo residual identificado y justificado (RT-02.11; SD4, sección 4.3.2.4); R8-05 lo trata con las medidas del SD4 —alarmas de retraso de replicación a los 5 y 15 min, reposición del enlace por el proveedor, preemisión de guías al cerrar la carga y conservación local en el NAS WORM de Talca— y la conmutación real previa al H5 y al H10 mide el RPO y el RTO.

No se autoriza corte sin continuidad medida ni aceptación sin las seis condiciones simultáneas del Art. 17.3 y acta según Art. 18. La suspensión del proveedor de lácteos, de marzo a septiembre de 2026, es antecedente ocurrido. La consulta V-13 del SD3 (Anexo 3.H) precisa en el mes 1, con el CLIENTE y el proveedor, qué evidencia de trazabilidad restablece la relación, y la Etapa 1 prioriza esa trazabilidad; la restitución depende del proveedor. La propuesta queda sujeta a las condiciones de cierre del Anexo 8.E, cada una con responsable, hito límite y evidencia; ninguna se da por cumplida sin las decisiones, pruebas y actas que allí se indican.

\section*{Referencias}
\addcontentsline{toc}{section}{Referencias}

Las fuentes citadas en este documento se listan a continuación en formato APA 7.ª edición.

\begin{itemize}[leftmargin=*,itemsep=3pt]
\item Distribuidora Puelche S.A. (2026a). \emph{Bases Administrativas de Licitación N.º TFEP-01/2026}, artículos 17, 18 y 50.2, y Formularios T-16 y T-22.
\item Distribuidora Puelche S.A. (2026b). \emph{Bases Técnicas Transversales de Licitación N.º TFEP-01/2026}, RT-07.04, RT-07.07, RT-19.04, RT-21.06, RT-21.07 y RT-26.04.
\item Distribuidora Puelche S.A. (2026c). \emph{Caso 02: Logística}, capítulos 10 a 14 y requisitos específicos citados.
\item Distribuidora Puelche S.A. (2026d). \emph{Aclaraciones de la Licitación N.º TFEP-01/2026}, secciones 2, 4, 6, 7 y 11 (Capítulo 8).
\item LafroX. (2026). Subdocumentos 1 a 4, 6 y 7, con los anexos y formularios citados.
\item International Electrotechnical Commission. (2018). \emph{IEC 60812:2018 Failure modes and effects analysis (FMEA and FMECA)}. IEC.
\item International Organization for Standardization. (2018). \emph{ISO 31000:2018 Risk management — Guidelines}. ISO.
\item Project Management Institute. (2017). \emph{La guía de los fundamentos para la dirección de proyectos (Guía del PMBOK®)} (6.ª ed.), capítulo 11. Project Management Institute.
\end{itemize}

\section*{Declaración de uso de IA}
\addcontentsline{toc}{section}{Declaración de uso de IA}

En cumplimiento de la sección 7.2 de las Aclaraciones de la licitación, la tabla siguiente declara el uso de herramientas de inteligencia artificial en este subdocumento, con la revisión humana de cada parte. La declaración se consolida en el Formulario A-6.

\renewcommand{\thetable}{8.IA.7}
\begin{tablalafrox}{Declaración de uso de IA}{tab:sd8-cuerpo-7}{>{\hsize=1.05136\hsize\linewidth=\hsize\RaggedRight\arraybackslash}X >{\hsize=0.77022\hsize\linewidth=\hsize\RaggedRight\arraybackslash}X >{\hsize=1.98104\hsize\linewidth=\hsize\RaggedRight\arraybackslash}X >{\hsize=0.46694\hsize\linewidth=\hsize\RaggedRight\arraybackslash}X >{\hsize=0.79930\hsize\linewidth=\hsize\RaggedRight\arraybackslash}X >{\hsize=0.93115\hsize\linewidth=\hsize\RaggedRight\arraybackslash}X}{\cab{Sección} & \cab{Herramienta} & \cab{Finalidad del uso} & \cab{Nivel en texto} & \cab{Nivel en diagramas} & \cab{Revisión humana (quién y qué verificó)}}
Introducción & Codex & Redacción a partir de SD2–SD4, SD6 y SD7 & Alto & Ninguno & [[REVISIÓN HUMANA]] \\
8.1 Plan de riesgos & Codex; Claude Code & Método, roles y escalas; Comité de Operación y textos de sección & Alto & Ninguno & [[REVISIÓN HUMANA]] \\
8.2 Identificación y Análisis de Riesgos & Codex; Claude Code & RBS textual, FMEA y escenarios; valor esperado en HH y simulación de Monte Carlo sobre el cronograma por actividad & Alto & Ninguno & [[REVISIÓN HUMANA]] \\
8.3 Plan de Acción a Riesgos & Codex; Claude Code & Respuestas, reservas y factibilidad; reserva de contingencia por valor esperado y su reparto por período; estrategias PMBOK, residual, riesgos secundarios y retorno de los controles & Alto & Ninguno & [[REVISIÓN HUMANA]] \\
Anexos 8.A a 8.F & Codex; Claude Code & Ver la declaración de los anexos & Alto & Ninguno & [[REVISIÓN HUMANA]] \\
Formulario T-16 & Codex; Claude Code & Ver la declaración del formulario & Alto & Ninguno & [[REVISIÓN HUMANA]] \\
Introducción, 8.1.1–8.1.3, 8.2.1–8.2.3, 8.3.2 y 8.3.3: correcciones de coherencia & Claude Code & Figura 8.1 (RBS) desde las fichas 8.A, contingencia adicional elegible, hitos no simulados, RPO residual alineado con el SD4, probabilidad del H9, soporte puente, fecha de la suspensión láctea, apetito de riesgo, mapa de los riesgos del Caso 19, costo-beneficio por riesgo, citas y referencias & Medio & Medio (Figura 8.1 en Mermaid a partir de las fichas del Anexo 8.A) & [[REVISIÓN HUMANA]] \\
\end{tablalafrox}
